% !TEX root = 06.tex
\documentclass[a4paper,10pt]{article}

\usepackage{pdfsync}
\usepackage[utf8]{inputenc}
\usepackage{microtype}
\usepackage{mathtools}
\usepackage{amsthm,amsfonts,amsmath,amssymb}
\usepackage{mathabx}
% \usepackage{MnSymbol}
\usepackage{hyperref}
\usepackage{bbold}
\usepackage[textwidth=2cm,textsize=small]{todonotes}
\usepackage{subcaption}
\usepackage{enumerate}
\usepackage{xspace}
\usepackage{stmaryrd}
\usepackage{verbatim}
\usepackage{cleveref}

\hypersetup{
    colorlinks,
    linkcolor={red!50!black},
    citecolor={blue!50!black},
    urlcolor={blue!80!black}
}

\author{}

\newcommand{\ignore}[1]{}

\newcommand{\st}{s.t.}
\newcommand{\ie}{i.e.}
\newcommand{\eg}{e.g.}

\newcommand{\sep}{\;|\;}
\renewcommand{\rule}[2]{\dfrac{#1}{#2}}
\newcommand{\set}[1]{\left\{#1\right\}}
\newcommand{\setof}[2]{\set{#1 \;\middle|\; #2}}
\newcommand{\tuple}[1]{\left( #1 \right)}
\newcommand{\pair}[2]{\tuple {#1, #2}}
\newcommand{\closure}[3]{\tuple {#1, #2, #3}}
\newcommand{\sem}[1]{\llbracket{#1}\rrbracket}
\newcommand{\powerset}[1]{\mathcal P(#1)}
\newcommand{\map}[2]{\mathsf{map}\; #1\; #2}
\newcommand{\filter}[2]{\mathsf{filter}\; #1\; #2}
\newcommand{\floor}[1]{\left\lfloor #1 \right\rfloor}
\newcommand{\PreRel}[2]{\mathsf{Pre}_{#1}(#2)}
\newcommand{\Pre}[1]{\PreRel {} {#1}}
\newcommand{\PreStar}[1]{\mathsf{Pre}^*(#1)}
\newcommand{\Post}[1]{\mathsf{Post}(#1)}
\newcommand{\PostStar}[1]{\mathsf{Post^*}(#1)}
\newcommand{\Conf}{\mathsf{Conf}}
\newcommand{\upwardclosure}[1]{#1\uparrow}
\newcommand{\card}[1]{\left| #1 \right|}
\newcommand{\lang}[1]{L(#1)}
\newcommand{\goesto}[1]{\xrightarrow{#1}}
\newcommand{\e}{\varepsilon}

\newcommand{\N}{\mathbb N}
\newcommand{\Z}{\mathbb Z}
\newcommand{\Q}{\mathbb Q}

% LTL
\newcommand{\true}{\mathbf{true}}
\newcommand{\false}{\mathbf{false}}
\newcommand{\X}{\mathbf X}
\newcommand{\U}{\mathbf U}
\newcommand{\Release}{\mathbf R}
\newcommand{\F}{\mathbf F}
\newcommand{\G}{\mathbf G}

% CTL
\newcommand{\EX}{\mathbf{EX}}
\newcommand{\AX}{\mathbf{AX}}
\newcommand{\EF}{\mathbf{EF}}
\newcommand{\AF}{\mathbf{AF}}
\newcommand{\EG}{\mathbf{EG}}
\newcommand{\AG}{\mathbf{AG}}
\newcommand{\EU}[2]{\mathbf E(#1 \U #2)}
\newcommand{\AU}[2]{\mathbf A(#1 \U #2)}
\newcommand{\ER}[2]{\mathbf E(#1 \Release #2)}
\newcommand{\AR}[2]{\mathbf A(#1 \Release #2)}

\newcommand{\NL}{\textsf{NL}}
\newcommand{\PTIME}{\textsf{PTIME}}
\newcommand{\PSPACE}{\textsf{PSPACE}}

\makeatletter

\def\dasharrowfill@#1#2#3#4{%
        $\m@th
        \thickmuskip0mu
        \medmuskip\thickmuskip
        \thinmuskip\thickmuskip
        \relax
        #4#1\mkern2mu
        \xleaders\hbox{$#4\mkern2mu#2\mkern2mu$}\hfill
        \mkern2mu
        #3$%
}

\def\dashleftarrowfill@{\dasharrowfill@\leftarrow\relbar\relbar}
\def\dashrightarrowfill@{\dasharrowfill@\relbar\relbar\rightarrow}
\def\dashleftrightarrowfill@{\dasharrowfill@\leftarrow\relbar\rightarrow}
\def\dashLeftarrowfill@{\dasharrowfill@\Leftarrow\Relbar\Relbar}
\def\dashRightarrowfill@{\dasharrowfill@\Relbar\Relbar\Rightarrow}
\def\dashLeftrightarrowfill@{\dasharrowfill@\Leftarrow\Relbar\Rightarrow}

\providecommand*\xdashleftarrow[2][]{%
  \ext@arrow 0055{\dashleftarrowfill@}{#1}{#2}}
\providecommand*\xdashrightarrow[2][]{%
  \ext@arrow 0055{\dashrightarrowfill@}{#1}{#2}}
\providecommand*\xdashleftrightarrow[2][]{%
  \ext@arrow 0055{\dashleftrightarrowfill@}{#1}{#2}}
\providecommand*\xdashLeftarrow[2][]{%
  \ext@arrow 0055{\dashLeftarrowfill@}{#1}{#2}}
\providecommand*\xdashRightarrow[2][]{%
  \ext@arrow 0055{\dashRightarrowfill@}{#1}{#2}}
\providecommand*\xdashLeftrightarrow[2][]{%
  \ext@arrow 0055{\dashLeftrightarrowfill@}{#1}{#2}}

\def\slashedarrowfill@#1#2#3#4#5{%
$\m@th\thickmuskip0mu\medmuskip\thickmuskip\thinmuskip\thickmuskip
  \relax#5#1\mkern-7mu%
  \cleaders\hbox{$#5\mkern-2mu#2\mkern-2mu$}\hfill
  \mathclap{#3}\mathclap{#2}%
  \cleaders\hbox{$#5\mkern-2mu#2\mkern-2mu$}\hfill
  \mkern-7mu#4$%
}

\def\rightslashedarrowfill@{%
  \slashedarrowfill@\relbar\relbar{\raisebox{1.2pt}{$\scriptscriptstyle\diagup$}}\rightarrow}
\newcommand\xslashedrightarrow[2][]{%
  \ext@arrow 0055{\rightslashedarrowfill@}{#1}{#2}}
\makeatother

\newtheorem{exercise}{Exercise}

\newenvironment{solution}{\begin{proof}[\protect{\textnormal{\emph{Solution.}}}]}{\end{proof}}

\let\oldqedhere\qedhere

\newif\ifstartedinmathmode
\renewcommand*{\qedhere}{%
  \relax\ifmmode\startedinmathmodetrue\else\startedinmathmodefalse\fi%
  \ifstartedinmathmode\tag*{\protect\oldqedhere}\else\ifinlist{\hfill\oldqedhere}{\oldqedhere}\fi%
}

\crefname{equation}{}{}
\crefname{exercise}{Exercise}{Exercises}


\title{Formal verification - Tutorial 06}
\date{23/03/2026}

\begin{document}

\maketitle

\section*{Timed automata}

\begin{exercise}
    Show that timed regular languages are effectively closed under
    \begin{enumerate}
        \item union,
        \item intersection,
        \item projection.
    \end{enumerate}
\end{exercise}
\begin{solution}
    For union and intersection, we can perform standard product constructions.
    For clocks we take the disjoin union of clocks.
    For projection, just project the alphabet.
\end{solution}

\begin{exercise}[Diagonal constraints]
    A \emph{diagonal constraint} is a constraint of the form
    %
    \begin{align*}
        x - y \sim c,
    \end{align*}
    %
    for some clock variables $x, y \in X$,
    comparison operator ${\sim} \in \set{<, \leq, =, \geq, >}$
    and integer constant $c \in \Z$.
    %
    Show that adding diagonal constraints to the syntax of timed automata
    does not increase their expressiveness
    (as recognisers of timed languages).
\end{exercise}
\begin{solution}
    The main intuition is that $x - y$ is invariant under time elapse.
    Thus we can store in the state whether $x - y \sim c$ holds or not.
    %
    When $y$ is reset, we need to know whether $x \sim c$ holds in order to update the state.
    We can add additional such tests to the guard of the transitions.
    %
    When $x$ is reset, the reasoning is similar.
    The construction is exponential in the number of diagonal constraints to be removed.
\end{solution}

\begin{exercise}
    Fix a set of clocks $X = \set{x_1, \dots, x_k}$ and a maximal constant $M \in \N$.
    Find an upper bound on the number of regions.
\end{exercise}
\begin{solution}
    A region is completely determined by:
    \begin{enumerate}
        \item The integral part of each clock up to $M$, or unbounded.
        There are $M + 2$ possibilities for each clock, thus $(M + 2)^k$ possibilities in total.
        \item The total order of the fractional parts of the clocks, of which there are $k!$ possibilities.
        \item For each clock, whether its fractional part is zero or not, which gives $2^k$ possibilities.
    \end{enumerate}
    This gives the upper bound
    %
    \begin{align*}
        (M + 2)^k \cdot k! \cdot 2^k. \qedhere
    \end{align*}
\end{solution}

\begin{exercise}
    Show that the nonemptiness problem for timed automata is \PSPACE-hard.
\end{exercise}
\begin{solution}
    We reduce from the nonemptiness problem for the intersection of $k$ regular languages,
    which is \PSPACE-hard (in fact \PSPACE-complete).
    %
    The idea is that a discrete-time clock can store the current state of a finite automaton.
    State updates are simulated by suitable time elapses and clock resets.

    There are many ways to achieve this.
    To fix ideas, consider finite automata $A_1, \dots, A_k$ with pairwise disjoint sets of states $Q_1, \dots, Q_k$,
    with a common set of states $Q = Q_1 \cup \cdots \cup Q_k = \set{q_1, \dots, q_m}$.
    %
    We use clocks $x_1, \dots, x_m$ to store the current state of the finite automata.
    The idea is that $x_i$ is zero if the relevant automaton is not in $q_i$
    and that $x_i$ is one if the relevant automaton is in $q_i$.
    %
    We need to efficiently update this information.
    %
    More generally, we show how we can set $x_i$ to a new value in $\set{0, 1}$,
    while keeping the other clocks unchanged.
    %
    Setting $x_i$ to zero is easy: Just reset it.
    %
    Setting $x_i$ to one (while keeping the other clocks unchanged) is more tricky.
    %
    We show how to do this with an example.
    %
    Suppose we have three clocks $x_1, x_2, x_3$, starting at $\tuple{0, 0, 1}$,
    and we want to set $x_1$ to one while keeping $x_2$ and $x_3$ unchanged,
    as to reach $\tuple{1, 0, 1}$.
    %
    More generally, we want to go from $\tuple{0, x_2, x_3}$ to $\tuple{1, x_2, x_3}$.
    %
    \begin{equation*}
        \begin{tabular}{c|c|c|c|c|c}
                & I & II & III & IV & V\\
        \hline
        $x_1$   & 0 & 1  & 0   & 1  & 1 \\
        $x_2$   & 0 & 1  & 1   & 2  & 0 \\
        $x_3$   & 1 & 2  & 0   & 1  & 1
        \end{tabular}
    \end{equation*}
    %
    In the first step, we elapse one time unit.
    In the second step, we reset $x_1$ and all clocks with value $2$.
    In the third step, we elapse one time unit.
    In the fourth step, we reset all clocks with value $2$.

    We can also implement a similar simulation using dense-time clocks,
    by encoding the state of the finite automaton in the total ordering of the fractional part of the clock values.
\end{solution}

\begin{exercise}
    Consider timed automata over a single clock $X = \set x$.
    What is the computational complexity of the nonemptiness problem?
\end{exercise}
\begin{solution}
    Regions for a single clock are considerably simpler.
    In fact, it suffices to keep track of the following information:
    %
    \begin{enumerate}
        \item The integral part of the clock up to $M$, or unbounded.
        \item Whether the fractional part of the clock is zero or not.
    \end{enumerate}
    %
    When clock constraints are encoded in binary, $M$ is still exponential in the size of the input.
    However, let $c_1 < \cdots < c_k$ be the pairwise distinct constants appearing in the input.
    It suffices just to keep track of whether the integral part of the clock is equal to $c_i$ for some $i$,
    or whether it is between $c_i$ and $c_{i + 1}$ for some $i$, or whether it is greater than $c_k$.
    %
    This gives a polynomial number of regions.
    Thus the nonemptiness problem for timed automata with a single clock is in \NL.
    Reachability is \NL-hard already for finite graphs, so altogether the problem is \NL-complete.
\end{solution}

\begin{exercise}
    What is the computational complexity of the universality problem for deterministic timed automata?
\end{exercise}
\begin{solution}
    Deterministic timed automata are efficiently closed under complementation,
    so the universality problem reduces to the nonemptiness problem in \PTIME,
    and thus it is in \PSPACE.
    %
    The problem is also \PSPACE-hard, since the nonemptiness problem for nondeterministic timed automata
    is \PSPACE-hard already for deterministic timed automata,
    and the latter reduces to the universality problem via complementation.
\end{solution}

\begin{exercise}
    Show that the universality problem for nondeterminstic timed automata with two clocks is undecidable.
\end{exercise}
\begin{solution}
    We reduce from the reachability problem for Turing machines, which is undecidable.
    %
    More precisely, we show that the \emph{complement} of the set of correct encodings of computations of a Turing machine is timed regular.
    %
    The idea is to encode a configuration of a Turing machine
    as a timed word lasting one time unit.
    %
    Two equal consecutive configurations are encoded by requiring that corresponding symbols appear at one time unit distance.
    %
    A timed automaton can check mistakes in this encoding using only two clocks.
    %
    A clock is used to keep track of the current time withing a configuration.
    The second clock is used to check that corresponding symbol in the next configuration is incorrect.
\end{solution}

\bibliographystyle{plainurl}
\bibliography{bibliography}

\end{document}